Nominally identical vibration-driven robots differ in speed, turning
bias and dead zone, and the differences drift with battery and wear. We
study the calibration of GRASPion, an open two-motor brainbot, as a
problem in its own right: what is the simplest model of its planar
motion from its two motor powers, how different are nominally identical
robots and how fast does one drift, and how few measurements does a new
robot need when a few calibrated robots exist. We present a grey-box
forward map (a per-motor dead zone and power law feeding a twist linear
in the two activations), show that steady-state data identify only the
mean of the clockwise and counter-clockwise centre-of-rotation offsets
that the robot's published digital twin treats as separate parameters,
derive a closed-form inverse with an explicit feasible set, and
formulate per-robot calibration as maximum-a-posteriori estimation under
a population prior with segments chosen by Bayesian D-optimal design.
On the lab's published 7×7 tables the map matches low-order polynomials
in leave-one-cell-out error (0.50 cm/s, 0.23 rad/s). On a synthetic
population the pipeline recovers the injected variability and reduces a
new robot's error with each added segment; those numbers are not
measurements. \hw{the campaign on N robots over D days: variance
decomposition, error against segments, spread under a common command}
